\section{Introduction}
\label{sec:intro}

Labels are expensive and images are cheap. Contrastive learning offers two ways to use a few labels together
with many images. Self-supervised methods such as \simclr \citep{chen2020simclr} pretrain on every image
without labels, by pulling two augmented views of the same image together; the labels are used only
afterwards, to train a classifier. Supervised contrastive learning, \supcon \citep{khosla2020supcon}, puts the
labels inside the pretraining loss: every image of the same class becomes a positive. With full labels \supcon
beats \simclr \citep{khosla2020supcon}, and \simclr itself approximately optimises a \supcon-like objective
\citep{luthra2025approximately}. It is therefore natural to assume that when \emph{some} labels exist, putting
them into the contrastive loss is the better choice.

This paper started from a result that seemed to contradict this assumption. In an earlier study on CIFAR-10,
we trained a ResNet-50 with \supcon on the 10\% of images that carry labels and a ResNet-50 with \simclr on
all images, each for ``100--200 epochs''. \supcon reached 75.2\% linear-probe accuracy, 9.1 points below
label-free \simclr (84.3\%; \cref{app:report}). We read this as a cost of lost diversity: restricting
pretraining to 5,000 labeled images seemed to destroy more than the labels could give back. A reviewer of
that report asked the questions this paper answers: \emph{Is the drop real or an optimisation artefact? When
does adding limited labels to a contrastive objective hurt, and why? And how does the result compare with the
methods the community actually uses with 10\% labels?}

\paragraph{What we did.}
We rebuilt the study around a controlled design (\cref{sec:setup}). All contrastive methods share one loss
(\cref{eq:loss}), one architecture, one augmentation pipeline, one evaluation code and, crucially, the same
number of gradient steps. A $2{\times}2$ design separates the two things that label-scarce pretraining changes
at once: whether the labels are used in the loss, and which images are in the pretraining pool. We evaluate
every encoder twice, with the standard linear probe that uses all 50,000 labels and with the 5,000 labels that
actually exist in the semi-supervised setting. Key runs are repeated over \nSeeds seeds and every key
difference is tested with a paired bootstrap; we add self-supervised (BYOL, SimSiam), supervised and
semi-supervised (FixMatch, MixMatch) baselines on the same split, and repeat the main comparison on CIFAR-100
and with a ResNet-50.

\paragraph{What we found.}
\begin{enumerate}[leftmargin=*]
\item \textbf{The original gap was an epoch-counting artefact} (\cref{sec:trap}). ``100 epochs'' over 5,000
images is ten times fewer gradient steps than 100 epochs over 50,000; in our report the difference was
21-fold. A ResNet-50 \supcon given the report's 970 steps reaches \rFiftySupconTen in our pipeline, close to
the 75.2\% we reported.
\item \textbf{Labels buy speed, unlabeled images buy accuracy} (\cref{sec:trap}). At equal steps, label-only
\supcon is far ahead early (\supconTenLin vs.\ \simclrTenLin after 970 steps), stops improving at about 84\%,
and is overtaken by \simclr, which reaches \simclrTwoHundredLin after 19,400 steps (\gapLong points ahead).
\item \textbf{Labels used in pretraining are spent} (\cref{sec:decompose}). With only the 5,000 labels available
after pretraining, fine-tuning a \simclr encoder gives \simclrTwoHundredFt, while encoders that used the same
labels inside the loss reach \supconTwoHundredFt (\supcon) and \unionTwoHundredFt (\supcon plus all unlabeled
images). The standard all-label probe hides this and even ranks the label-using encoder first at shorter
budgets.
\item \textbf{The mechanism is memorisation, and coverage limits it} (\cref{sec:geometry}). Label-only \supcon
collapses the classes of the images it has seen \supconCdnvRatio times more tightly than those of new images,
so the labeled images are already fitted when the labels are used again. The reviewer's proposed diversity
metric, the spread of augmented views within a class, is the same for a random 1\% subset (\spreadOne) and the
full data set (\spreadFull); what a small pool lacks is coverage of unseen images, and pools made redundant at a
fixed size give worse encoders.
\end{enumerate}
On the reviewer's last question, the answer is sobering for contrastive pretraining: MixMatch and FixMatch
\citep{berthelot2019mixmatch,sohn2020fixmatch}, given the same image budget, beat every contrastive pipeline we
trained (\mixmatchAcc and \fixmatchAcc; \cref{sec:baselines}). This is also consistent with our explanation:
both methods use the labels to create targets for the unlabeled images instead of clustering the labeled ones. We also quantify two
observations of the original study, the colour shortcut behind the large effect of colour
jitter (\cref{sec:shortcut}) and a resolution failure on ImageNet photographs (\cref{sec:ood}).

\paragraph{Why this matters.}
The practical message is short: \emph{pretrain on every image without labels, and spend the labels last}.
Labels inside a contrastive loss make training faster, which matters when compute is tight; but with a
realistic budget, the same labels are worth more when they are kept for fine-tuning. The broader message is
about evaluation. A plausible mechanism (``fewer images means less diversity'') was attached to a 9-point gap
that was caused by a 21-fold difference in training; and the standard linear probe, which assumes all labels
are available, can rank methods in the opposite order from the protocol that matches the label budget. Both
echo earlier warnings about evaluating semi-supervised learning \citep{oliver2018realistic}.
